\section*{Capítulo \ref{ch:b3:organometallic-bonding} --- Química organometálica: ligação e ligantes}

\begin{solution}{exo:b3:organometallic-bonding:1}
$\ce{Fe(CO)5}$: $8 + 10 = 18$. $\ce{Mn2(CO)10}$: $7 + 10 + 1 = 18$ por Mn. $\ce{CpMo(CO)3H}$: $6 + 5 + 6 + 1 =
18$. $\ce{[PtCl3(C2H4)]-}$: $10 + 3 + 2 + 1 = 16$. $\ce{Cr($\eta^6$-C6H6)(CO)3}$: $6 + 6 + 6 = 18$. $\ce{Cp2ZrCl2}$:
$4 + 10 + 2 = 16$.
\end{solution}

\begin{solution}{exo:b3:organometallic-bonding:2}
Fe(0) $d^8$; Mn(0) $d^7$; Mo(II) $d^4$; Pt(II) $d^8$; Cr(0) $d^6$; Zr(IV) $d^0$.
\end{solution}

\begin{solution}{exo:b3:organometallic-bonding:3}
$\ce{[Mn(CO)6]+} > \ce{Cr(CO)6} > \ce{[V(CO)6]-}$: quanto mais rico em elétrons o metal, mais ele
retrodoa para $\pi^*$ e mais fracas as ligações C--O.
\end{solution}

\begin{solution}{exo:b3:organometallic-bonding:4}
Etano ($\ce{Mn(CO)5}$ e $\ce{CpFe(CO)2}$ são \omterm{def:b3:organometallic-bonding:isolobal}{isolobais} com \ce{CH3}); de novo etano; tetraedrano
$\ce{(CH)4}$ com três CH substituídos por $\ce{Co(CO)3}$.
\end{solution}

\begin{solution}{exo:b3:organometallic-bonding:5}
$fac$ ($C_{3v}$): duas (A$_1$ + E). $mer$ ($C_{2v}$): três (2A$_1$ + B$_1$).
\end{solution}

\begin{solution}{exo:b3:organometallic-bonding:6}
A fosfina volumosa congestiona o metal: a dissociação alivia a tensão, logo é
mais rápida, e os intermediários de baixa coordenação são favorecidos.
\end{solution}

\begin{solution}{exo:b3:organometallic-bonding:7}
O complexo de crômio é um \omterm{def:b3:organometallic-bonding:carbene}{carbeno de Fischer} (substituinte OMe, Cr(0), ligantes CO): uma
amina ataca o carbono carbênico eletrofílico e substitui o OMe. O alquilideno de
tântalo é um \omterm{def:b3:organometallic-bonding:carbene}{carbeno de Schrock}: seu carbono nucleofílico ataca a carbonila de uma
cetona, dando um alqueno e uma unidade Ta=O, como em uma reação de Wittig.
\end{solution}

\begin{solution}{exo:b3:organometallic-bonding:8}
Alternado: sem sobreposição $\delta$, os dois elétrons $\delta$ são não ligantes: ordem de ligação 3. $\ce{[Re2Cl8]^{4-}}$:
$\sigma^2\pi^4\delta^2\delta^{\ast2}$, ordem de ligação 3.
\end{solution}

\begin{solution}{exo:b3:organometallic-bonding:9}
Uma distância M$\cdots$H curta (difração, de preferência de nêutrons); um número de onda de estiramento C--H
baixo; um sinal de RMN em campo alto com um acoplamento C--H a uma ligação reduzido.
\end{solution}

\begin{solution}{exo:b3:organometallic-bonding:10}
16: $d^8$ quadrado-planar ($\ce{[PtCl4]^{2-}}$, catalisador de Wilkinson), cujo orbital do 18.º elétron
é alto demais; os \omterm{def:b3:organometallic-bonding:metallocene}{metalocenos} dos metais do início da série, como $\ce{Cp2ZrCl2}$, congestionados demais para mais
ligantes. 17: $\ce{V(CO)6}$, que não pode dimerizar por razões estéricas. 19: o cobaltoceno,
cujo elétron extra fica em um orbital antiligante (e se perde facilmente).
\end{solution}

\begin{solution}{exo:b3:organometallic-bonding:11}
O CO é um $\pi$-aceptor forte: ele abaixa o nível $t_{2g}$ e torna $\Delta_{\mathrm o}$ grande, muito acima da
energia de emparelhamento: $t_{2g}^6$, diamagnético. As transições $d$--$d$, em $\Delta_{\mathrm o}$, caem no ultravioleta,
e nenhuma luz visível é absorvida.
\end{solution}

\begin{solution}{exo:b3:organometallic-bonding:12}
A doação a partir de $\pi$ e a retrodoação para $\pi^*$ diminuem ambas a ordem de ligação C--C. No sal de Zeise
(Pt(II), retrodoação moderada), a ligação se alonga um pouco; um metal de baixa valência,
rico em elétrons, retrodoa fortemente e o complexo se aproxima do
limite do \omterm{def:b3:organometallic-bonding:dcd}{metalaciclopropano}, com uma ligação C--C próxima de uma ligação simples.
\end{solution}

\begin{solution}{pb:b3:organometallic-bonding:1}
\textbf{1.} $\ce{Fe^{2+}}$ ($d^6$) $+ 2 \times 6 = 18$.
\textbf{2.} $8 + 2 \times 5 = 18$.
\textbf{3.} $9 + 10 = 19$.
\textbf{4.} $10 + 10 = 20$.
\textbf{5.} $18 - 1 = 17$.
\textbf{6.} Fe(II) $d^6$; Co(II) $d^7$; Ni(II) $d^8$; Fe(III) $d^5$.
\textbf{7.} $a_{1g}$ ($d_{z^2}$), $e_{1g}$ ($d_{xz}$, $d_{yz}$), $e_{2g}$ ($d_{xy}$, $d_{x^2-y^2}$).
\textbf{8.} O par $e_{1g}$ se sobrepõe fortemente à combinação $e_{1g}$ preenchida dos anéis: orbitais
ligantes sobretudo nos anéis, antiligantes $e_{1g}^\ast$ sobretudo no metal, altos em
energia.
\textbf{9.} $e_{2g} \approx a_{1g} < e_{1g}^\ast$.
\textbf{10.} $e_{2g}^4a_{1g}^2$: nenhum elétron desemparelhado.
\textbf{11.} Um elétron em $e_{1g}^\ast$: um elétron desemparelhado, $1.73\,\mu_B$.
\textbf{12.} Dois elétrons nos dois orbitais $e_{1g}^\ast$, paralelos: dois desemparelhados, $\sqrt{8} =
2.83\,\mu_B$.
\textbf{13.} $e_{2g}^3a_{1g}^2$: um elétron desemparelhado. Um buraco $e_{2g}^3$ dá um estado orbitalmente degenerado
(E) com \omterm{def:b3:complex-spectra-magnetism:moment}{contribuição orbital}.
\textbf{14.} Seu 19.º elétron está em um orbital antiligante e é facilmente retirado,
dando o cátion cobaltocênio, de 18 elétrons.
\textbf{15.} O elétron retirado vem de um orbital quase não ligante: a
estrutura quase não muda (pequena \omterm{def:b3:complex-mechanisms:reorganisation}{energia de reorganização}, \cref{ch:b3:complex-mechanisms}).
\textbf{16.} O par é rápido e reversível na maioria dos solventes, o composto é
solúvel e estável, e seu potencial depende pouco do solvente, pois a
carga se espalha por um íon grande.
\textbf{17.} Com 20 elétrons, dois deles antiligantes, ele reage perdendo-os:
adiciona ligantes com deslizamento ou perda de um anel, ou é oxidado.
\textbf{18.} $\ce{CpFe(CO)2}$: $8 + 5 + 4 = 17$ elétrons, a um de 18, um orbital de fronteira: \omterm{def:b3:organometallic-bonding:isolobal}{isolobal}
com \ce{CH3}. O dímero $\ce{[CpFe(CO)2]2}$ é o ``etano'', com uma ligação Fe--Fe (na prática, dois ligantes CO
fazem ponte sobre ela).
\textbf{19.} $\ce{CpFe(CO)2-Mn(CO)5}$, com uma ligação Fe--Mn.
\textbf{20.} $\ce{PPh3}$ é melhor doador e pior $\pi$-aceptor que o CO: o metal retrodoa
mais para os CO restantes, cujos estiramentos diminuem.
\textbf{21.} Duas (A$_1$ + E).
\textbf{22.} O deslizamento do anel para $\eta^3$ libera o equivalente a dois elétrons de coordenação: um ligante
de entrada pode se ligar de modo associativo sem que o metal ultrapasse 18 elétrons.
\textbf{23.} \textbf{$\mu_{\text{spin puro}}(\text{niqueloceno}) = \sqrt{2 \times 4} = 2.83\,\mu_B$.}
\end{solution}
